\begin{multicols*}{2}

\ubsection{Introduction}{intro}

% {\color{Red} \ubsubsection{Lost traditions}{}}


{\color{Red} \ubsubsection{I. Historical Context}{historical-context}}
 % \& Problem Statement

Since the dawn of Christianity, Catholics have prayed the canonical hours and Psalms for millennia, continuing on the ancient Jewish Temple tradition, and our current time is no exception. From the height of the Middle ages with every order and secular cathedral singing its own Use of the Office, to laymen reciting various \textit{Libri Horarum}, with the rise of the printing press and the Tridentine standard, the number of local Uses has steadily decreased in popularity, leaving the Dominicans, along with a handful of other Sees and Orders hanging onto their Uses.

Catholics today wishing to pray the hours have a plethora of options, yet a strange dearth in the traditional variety. One can purchase the various editions of the \textit{nouveau} Liturgy of the Hours, to the various reprints of the 1962 secular and 1963 monastic breviaries, yet outside of those, we venture into the vernacular with editions of the same, with the added choice of various Eastern Offices, from Byzantine, to Eastern Syriac, to Coptic, to Armenian. Yet within the Roman Rite itself, we are left to a strange situation of being content with the new, reformed books of Vatican II, reprints of the 1962 books, or we are left outbidding each other on online auctions for ancient Breviaries last printed seventy or eighty years ago. The same holds even more true and dire for those wishing to bypass the Breviary reforms of Pius X.


{\color{Red} \ubsubsection{II. Vision \& Editorial Principles}{vision-and-editorial-principles}}

% II. The Project: Vision \& Editorial Principles

What I present here is a labor two years in the making, starting in June 2024 with these several goals:
\begin{itemize}
	\item To typeset a Pre-\textit{Divino Afflatu} Breviary,
	\item To recover the Dominican Rite, as standardized by Bl. Humbert of Romans,
	\item To show LaTeX is viable for typesetting Breviaries, providing an open-source template, and
	\item To demonstrate this using modern conveniences, such as Print-on-Demand services.
\end{itemize}

Thus, by the grace of God, the first \emph{Totum}, Dominican Rite, Latin Pre-55, and Pre-\textit{Divino Afflatu} Breviary printed in more than 60 years.

I have endeavoured to keeping the layout of this Breviary as faithfully as possible to the 13th-century manuscript of Humbert, namely, the Prototype of Santa Sabina XIV L1, which was the foundational text that unified early Dominicans in prayer.

This is the Rite that Ss. Thomas Aquinas, Albert the Great, and Raymond of Penyafort prayed, which has been restored and newly typeset, without the cluttered Kalendar that are characteristic of later editions.

By bypassing modern and Baroque accretions, this text allows contemporary laity, religious, and scholars to experience the Divine Office as it was laid down by Bl. Humbert, and as prayed by early Dominican saints. This Breviary bridges the gap between historical fidelity and accessibility for daily prayer.




{\color{Red} \ubsubsection{III. Content \& Structural Overview}{content-and-structual-overview}}


As a \emph{Totum} Breviary, this contains everything needed for a traveling Friar Preacher to recite daily, as well as the Little Office of the Blessed Virgin, along with rubrics on the communal obligation of the Office of the Dead, among other rubrics. I have included the Necrology of the Masters-Generals, as they are written in the Prototype, namely to 1345, and included the rubric on how to execute the termination of collects.

As I have kept faithful to the layout and order of the source manuscript, the rubrics are scattered in a variety of places, from the rubrics excerpted from the Collectarium (including the Kalendar), to the rubrics concerning the Temporal Cycle before Advent, to Rubrics concerning Sundays Per Annum placed after Epiphany, to rubrics concerning the Sanctoral placed before the Sanctoral Propers, placed in April are the Commons for Paschaltide, and rubrics at the very end governing the supplemental offices.

I have kept the medieval orthography, conceding to use modern u's and v's for clarity's sake.

It is recommended for the novice to acquire the Ordo and Kalendar, available as a companion to this Breviary, at \textit{https://breviariumpredicatorum.pages.dev/\#ordo} as a guide to the Rubrics and Kalendar for each year, with pointers to which Feast, Psalms, Antiphons, Responsories, and Memories are said on each day. It also includes an English version of the major rubrics in this Breviary.

For those interested in the History of the Dominican Rite, the excellent study, ``\emph{A History of the Dominican Liturgy}'' by Fr. William Bonniwell, O.P. is the standard, as well as the more recent history from 1945 to 1969 by the Very Rev. Fr. Augustine Thompson, O.P. is also available to read online.


{\color{Red} \ubsubsection{IV. Comparative Guide}{comparative-guide}}

The Dominican Breviary will be very much familiar to anyone who has recited the 1962 Roman Office on I Class feasts, with 7 canonical hours, and the nocturnal hours of Matins. As on I Class feasts, at the minor hours are recited Psalms 53 and 118, and at Compline, which is more or less fixed, are said the Sunday Psalms of 4, 90, and 133, with the addition of the first six verses of Psalm 30, and a rich variety of Antiphons.

At Matins, Sundays have 18 Psalms, ferias have 12, and feasts outside of Paschaltide have 9 Psalms, but in Paschaltide, 3 Psalms are said daily.

At Lauds 7 Psalms are said plus the Old Testament canticle, with the 3rd and 5th slot reserved daily for Psalms 62 \& 66, and Psalms 148-150 daily.

At Prime, outside of Sundays between Septuagesima and Palm Sunday, it is fixed to Psalms 53 and the first four stanzas of 118. On those excepted Sundays, Psalms 21-25 and 117 would be added. The Quicumque is recited every Sunday excepting when a Totum Duplex feast occurred. Those accustomed to the Roman Office will notice the Chapter Office of Prime is decoupled from Prime, and would instead be recited in choir after Lauds or Prime, depending on the Prior. After the Martyrology is read, the Lectio Brevis that follows is not that of the season or the Nones Chapter, but either the beginning of the Gospel of the day, if it is a Simplex feast (or above), or the beginning of the Constitutions of their Order. If a Friar did not have the proper Gospel of the day, the Prologue of John was read instead. The remaining minor hours are very much the same as the Roman Rite, finish the remainder of Psalm 118.

At Vespers, there are special rules compared to the Roman on which Psalms and antiphons were to be recited, depending on the rank of the feast, but on Sundays began Psalms 109-113 and ferias continuing at five per day, onto Psalm 147, excepting Psalms 117, 118, 133, and 142 which were recited elsewhere.


{\color{Red} \ubsubsection{V. Community \& Open-Source Collaboration}{community-and-open-source-collaboration}}


This Breviary is an open-source project. Despite rigorous proofreading against the source manuscript, typos other grammatical errors may persist in this first edition. If you find errors, or if you wish to contribute to the underlying LaTeX template, please submit an issue or pull request at: \emph{https://www.Github.com/AlanL399/BreviaryTypesetting/}

For the Order of Preachers, the solemn celebration of the Divine Office is their primary duty and the source of their contemplative and doctrinal life. May this edition serve to nourish your own contemplation, that you may bear its fruits to the world.


{\color{Red} \ubsubsection{VI. How to Navigate}{how-to-navigate}}
Since this Breviary was printed via Print-on-Demand, there are no ribbons, and users are encouraged to use plenty of prayer cards and sticky tabs to keep their place in this volume. The following are recommended, but not limited to:

Bookmarks/Prayer Cards
\begin{itemize}
	\item Psalms
	\item Temporal
	\item Sanctoral
	\item Common
	\item Common Lessons
\end{itemize}


Sticky Tabs
\begin{itemize}
	\item Canticles/Te Deum
	\item Hours (1st Sunday After Octave of the Epiphany)
	\item Blessings/Preciosa
	\item Daily Commemoration of St. Dominic
	\item Confiteor/Salve
	\item Compline Benedictio (Collectarium)
\end{itemize}

I have included two indices for the psalter and hymns in the back of the Breviary, and the Table of Contents lists all the sections that are used daily.


{\color{Red} \ubsubsection{Test}{test-foreward}}

\selectcolormodel{gray}

{\color{Red} \ubsubsection{Test 2}{test2-foreward}}



\textcolor{black!100}{the five boxing wizards jump quickly 100}
\par \textcolor{black!90}{the five boxing wizards jump quickly 90}
\par \textcolor{black!80}{the five boxing wizards jump quickly 80}
\par \textcolor{black!70}{the five boxing wizards jump quickly 70}
\par \textcolor{black!60}{the five boxing wizards jump quickly 60}
\par \textcolor{black!50}{the five boxing wizards jump quickly 50}


\lettrine[lines=2]{\zallmancaps \color{Red} R}{ubric} that has sections of text \ul{with underline} and \emph{italic text.} And more text. \ul{This is also some unique text that will be tested upon with the eyes, so don't be afraid to read on.} \emph{Here is some more text that is italicized, so please read on to test the reader. \ul{And we underline it as well, for fun's sake, to add more text to read.}} Red
\lettrine[lines=2]{\zallmancaps \color{Blue} R}{ubric} that has sections of text \ul{with underline} and \emph{italic text.} And more text. \ul{This is also some unique text that will be tested upon with the eyes, so don't be afraid to read on.} \emph{Here is some more text that is italicized, so please read on to test the reader. \ul{And we underline it as well, for fun's sake, to add more text to read.}} Blue


\textcolor{black!100}{waltz bad nymph for quick jigs vex 100}
\par \textcolor{black!90}{waltz bad nymph for quick jigs vex 90}
\par \textcolor{black!80}{waltz bad nymph for quick jigs vex 80}
\par \textcolor{black!70}{waltz bad nymph for quick jigs vex 70}
\par \textcolor{black!60}{waltz bad nymph for quick jigs vex 60}
\par \textcolor{black!50}{waltz bad nymph for quick jigs vex 50}

\begin{responsory-breve}
{Sample Responsory With Text * And A Second Part.}{Plus The Versicle.}
\end{responsory-breve}%

\newline Otherwise, there are no italicized text, \ul{only underlined text.} \emph{So, please, dear reader, make the right choice.}

\grecross \grealtcross \gothRbar \gothVbar \ \greheightstar \ \gresixstar \ 
\grecross \grealtcross \gothRbar \gothVbar \ \greheightstar \ \gresixstar

% {\color{Red} \ubsubsection{Test}{-foreward}}
% {\color{Red} \ubsubsection{Test}{-foreward}}
% {\color{Red} \ubsubsection{Test}{-foreward}}


% \lettrine[lines=2]{\zallmancaps \color{Red} O}{}
% \lettrine[lines=2]{\zallmancaps \color{Blue} I}{n}
% \lettrine[lines=2]{\zallmancaps \color{Red} O}{}
% \lettrine[lines=2]{\zallmancaps \color{Blue} I}{n}

\end{multicols*}